\documentclass[10pt]{article}

\usepackage[utf8]{inputenc}

\usepackage[T1]{fontenc}

\usepackage{amsmath}

\usepackage{amsfonts}

\usepackage{amssymb}

\usepackage[version=4]{mhchem}

\usepackage{stmaryrd}

\usepackage{graphicx}

\usepackage[export]{adjustbox}

\graphicspath{ {./images/} }



\begin{document}

Для У. м. принята стандартизация



\[

\begin{gathered}

P_{n}(x, \lambda) \equiv C_{n}^{(\lambda)}(x)= \\

=\frac{(-2)^{n}}{n^{\prime}} \cdot \frac{\Gamma(n+\lambda) \Gamma(n+2 \lambda)}{\Gamma(\lambda) \Gamma(2 n+2 \lambda)}\left(1-x^{2}\right)^{-\lambda+\frac{1}{2}} \frac{d^{n}}{d x^{n}} \times \\

\times\left[\left(1-x^{2}\right)^{\left.n+\lambda-\frac{1}{2}\right]}\right.

\end{gathered}

\]



и имеет место представление



\[

C_{n}^{(\lambda)}(x)=\sum_{k=0}^{\left[\frac{n}{2}\right]}(-1)^{k} \frac{\Gamma(n-k+\lambda)}{\Gamma(\lambda) k^{\prime}(n-2 k)^{\prime}}(2 x)^{n-2 k} .

\]



У. м. являются коэффициентами разложения в степенной ряд производящей функции



\[

\frac{1}{\left(1-2 x w+w^{2}\right)^{\lambda}}=\sum_{n=0}^{\infty} C_{n}^{(\lambda)}(x) w^{n} .

\]



У. м. $C_{n}^{(\lambda)}(x)$ удовлетворяет дифференциальному уравнению



\[

\left(1-x^{2}\right) y^{\prime \prime}-(2 \lambda+1) x y^{\prime}+n(n+2 \lambda) y=0 .

\]



Наиболее употребительны формулы



$(n+1) C_{n+1}^{(\lambda)}(x)=2(n+\lambda) x C_{n}^{(\lambda)}(x)-(n+2 \lambda-1) C_{n-1}^{(\lambda)}(x)$,



\[

C_{n}^{(\lambda)}(-x)=(-1)^{n} C_{n}^{(\lambda)}(x),

\][



\textbackslash frac\{d\}\{d x\}\textbackslash left[C\_\{n\}\^{}\{(\textbackslash lambda)\}(x)\textbackslash right]=2 \textbackslash lambda C\_\{n-1\}\^{}\{(\textbackslash lambda+1)\}(x),



][



C\_\{n\}\^{}\{(\textbackslash lambda)\}( \textbackslash pm 1)=( \textbackslash pm 1)\^{}\{n\} \textbackslash frac\{2 \textbackslash lambda(2 \textbackslash lambda+1) ···(2 \textbackslash lambda+n-1)\}\{n !\}=



][



=( \textbackslash pm 1)\^{}\{n\}\textbackslash left(\[

\begin{array}{c}

n+2 \lambda-1 \\

n

\end{array}

\]\textbackslash right) \textbackslash text \{. \}

]



Лит. см. при ст. Ортогональные многочлены.



УЈЬТРАФИЛЬТР - фильтр, являющийся максимальным в том смысле, что всякий содержащий его фильтр совпадает с ним.



У. можно определить как систему подмножеств, удовлетворяющую трем условиям: 1) пустое множество ей не принадлежит; 2) пересечение двух принадлежащих ей подмножеств также ей принадлежит; 3) для любого подмножества либо оно само, либо его дополнение принадлежит этой системе.



Все У. делятся на два класса: тривиальные (или фиксированные) и свободные. У. наз. т р и в и а л ь н ы м, если он представляет собой систему всех подмножеств, содержащих нек-рую точку, такой У. наз. также фиксированным на этой точке. У. наз. с в о б о д н ы м, если пересечение всех его элементов есть пустое множество, другими словами, если он не фиксирован ни на какой точке. Существование свободных У. недоказуемо без Выбора аксиомы.



Для каждого фильтра имеется содержащий его У., более того, каждый фильтр есть в точности пересечение всех содержащих его У.



Лит.. [1] Б у р б а к и Н, Общая топология. Основные



\begin{center}

\includegraphics[max width=\textwidth]{2023_07_09_c38c849daff22065fc4ag-1}

\end{center}



УМНОЖЕНИЕ ч и с е л - одна из основных арифметич. операций. У. заключается в сопоставлении двум числам $a$ и $b$ (вазываемым с о м в о ж и т е л я м и) третьего числа $c$ (называемого п р о и 3 в е д е н и е м). У. обозначается знаком $X$ или - ; в буквенном обозначении эти знаки, как правило, опускаются.



У. целых положительных чисел определяется следующцим образом через сложение: произведением чисел $a$ и $b$ считается число $c$, равное сумме $b$ слагаемых, каждое из к-рых равно $a$, так что



\[

a b=\underbrace{a+a+\ldots+a}_{\text {b раз }} .

\]



Число $a$ при этом наз. м н о ж и м м,$b-$ м н о ж ит е л ем. У. положительных рациональных чисел $\frac{m}{n}$ и $\frac{p}{q}$ определяется равенством



\[

\frac{m}{n} \cdot \frac{p}{q}=\frac{m p}{n q}

\]



(см. Дробъ). Произведение двух отрицательных сомножителей положительно, а положительного и отрицательного - отрицательно, причем модуль произведения в том и другом случае равен произведению модулей сомножителей. Произведение иррациональных чисел определяется как предел произведений их рациональных приближений. У. комплексных чисел $\alpha=a+b i$ и $\beta=$ $=c+d i$ задается формулой



\[

\alpha \beta=(a+b i)(c+d i)=a c-b d+(a d+b c) i

\]



или, в тригонометрич. форме, $\left(\alpha=r_{1}\left(\cos \varphi_{1}+i \sin \varphi_{1}\right)\right.$, $\left.\beta=r_{2}\left(\cos \varphi_{2}+i \sin \varphi_{2}\right)\right)$



\[

\alpha \beta=r_{1} r_{2}\left(\cos \left(\varphi_{1}+\varphi_{2}\right)+i \sin \left(\varphi_{1}+\varphi_{2}\right)\right) .

\]



У. чисел коммутативно, ассоциативно и дистрибутивно слева и справа относительно сложения (см. Коммутативность, Ассоциативность, Дистрибутивность). При этом $a \cdot 0=0, \quad a \cdot 1=a$.



В общей алгебре У. может называться любая алгебрачческая операчия ( $n$-арная, $n \geqslant 2)$; чаще всего бинарвая операция (әруппоид). В нек-рых случаях эти операции являются обобщением обычного У. чисел. Напр., У. кватернионов, У. матриц, У. преобразований. Однако свойства У. чисел (напр., коммутативность) в этих случаях могут утрачиваться.



О. А. Иванова. УНАРНАЯ АЛГЕБРА, у н о и д,- универсальная алгебра $<A,\left\{f_{i} \mid i \in I\right\}>$ с семейством $\left\{f_{i} \mid i \in I\right\}$ унарных операций $f_{i}: A \rightarrow A$.



Важный пример У. а. дает групповой гомоморфизм $\varphi: G \rightarrow S_{A}$ произвольной групшы $G$ в группу $S_{A}$ всех подстановок множества $A$. Такой гомоморфизм наз. д е й с т в и е м группы $G$ на $A$. Определяя унарную операцию $f_{g}: A \rightarrow A$ для каждого элемента $g \in G$ как подстановку $\varphi(g)$ из $S_{A}$, отвечающую элементу $g$ при гомоморфизме $\varphi$, получают У. а. $<A,\left\{f_{g} \mid g \in G\right\}>$, в к-рой



$f_{1}(x)=x, f_{g}\left(f_{h}(x)\right)=f_{g h}(x), \quad x \in A, g, h \in G$.



Структуру У. а. несет на себе любой модуль над кольцом. Каждый детерминированный полуавтомат с множеством состояний $S$ и входными символами $a_{1}, \ldots$, $a_{n}$ также можно рассматривать как У. а. $\left\langle S, f_{1}, \ldots\right.$, $f_{n}>$, в к-рой $f_{i}(s)=a_{i} s$ есть состояние, следующее за состоянием $s$ в зависимости от входного символа $a_{i}$.



У. а. с одной основной операцией наз. м о н о у в а рн о й, или у н а $\mathrm{p}$ о м. Примером унара может служить алгебра Пеано $\langle P, f>$, где $P=\{1,2, \ldots\}$ и $f(n)=$ $=n+1$.



Тождества произвольной У. а. могут быть лишь следующцих типов:



\[

\begin{gathered}

\mathrm{I}_{1} \cdot f_{i_{1}} \ldots f_{i_{k}}(x)=f_{j_{1}} \ldots f_{i_{l}}(x), \\

\mathrm{II}_{1} \cdot f_{i_{1}} \ldots f_{i_{k}}(x)=f_{j_{1}} \ldots f_{i_{l}}(y), \\

\mathrm{I}_{2} \cdot f_{i_{1}} \ldots f_{i_{k}}(x)=x, \quad \mathrm{II}_{2} \cdot f_{i_{1}} \ldots f_{i_{k}}(x)=y, \\

\mathrm{I}_{3} . x=x, \quad \mathrm{II}_{3} . x=y .

\end{gathered}

\]



Тождество $\mathrm{II}_{2}$ равносильно тождеству $\mathrm{II}_{3}$, выполнимому лишь в одноэлементной алгебре. Многообразие У. а., определяемое лишь тождествами вида $\mathrm{I}_{1}, \mathrm{I}_{2}$ или $\mathrm{I}_{3}$, наз. р е г у ля р н о о п р е де ли мы м. Существует следующая связь между р е г у ля р н о оп р е д ел и м ы и многообразиями У. а. и полугрупшами (cM. [1], [3], [4]).





\end{document}